\documentclass[11pt]{article}
\usepackage[margin=1in]{geometry}
\usepackage[T1]{fontenc}
\usepackage[utf8]{inputenc}
\usepackage{lmodern}
\usepackage{amsmath,amssymb,mathtools,booktabs,hyperref}

\title{Penalized Trend Estimation from First Principles: A Step-by-Step Companion}
\author{Heriberto Espino Montelongo}
\date{\today}

\newcommand{\R}{\mathbb{R}}

\begin{document}
\maketitle

\begin{abstract}
This companion document explains the mathematical and computational ideas used
in the Trend Estimation library for a reader who has not previously studied
penalized smoothing. It develops finite differences, quadratic penalties,
closed-form trend estimation, temporal forecasting, analytic sensitivity, and
numerical smoothness selection. The active research paper is separate; this
document is pedagogical and must remain consistent with the tested library and
the derivations recorded under \texttt{notes/}.
\end{abstract}

\section{The pure penalized smoother}
Let $y=(y_1,\ldots,y_n)^\top\in\R^n$ and let $D_d$ compute finite
differences of order $d$. We define
\[
Q=D_d^\top D_d.
\]
The pure penalized trend solves
\[
\min_t\;\|y-t\|_2^2+\lambda\|D_dt\|_2^2,
\]
so
\[
\boxed{
\widehat t_\lambda=(I+\lambda Q)^{-1}y.
}
\]
The implementation solves the corresponding linear/spectral system rather than
forming an explicit dense inverse.

\section{Sensitivity to the penalty}
Writing
\[
S_\lambda=(I+\lambda Q)^{-1},
\]
the inverse-matrix derivative gives
\[
S_\lambda'=-S_\lambda Q S_\lambda.
\]
Therefore
\[
\boxed{
\widehat t_\lambda'=-S_\lambda Q\widehat t_\lambda,
}
\]
and
\[
\boxed{
\widehat t_\lambda''
=
2S_\lambda Q S_\lambda Q\widehat t_\lambda.
}
\]

\section{Forecast loss}
At a forecast origin $T$, no future observation may influence the fitted trend.
Let $H$ denote the continuation operator and let
\[
r_T(\lambda)
=
y_{T+1:T+h}
-
H S_\lambda y_{\mathrm{past}}.
\]
Then
\[
r_T'
=
H S_\lambda Q S_\lambda y_{\mathrm{past}},
\]
and for
\[
f_T(\lambda)=\frac1h r_T^\top r_T,
\]
we obtain
\[
\boxed{
f_T'(\lambda)
=
\frac{2}{h}
r_T^\top
H S_\lambda Q S_\lambda y_{\mathrm{past}}.
}
\]
A second derivative is available as well and is derived step by step in
\texttt{notes/derivative.md}.

\section{Why optimize log-lambda?}
Set
\[
\theta=\log\lambda,
\qquad
\lambda=e^\theta.
\]
If $g(\theta)=f(e^\theta)$, then
\[
g'(\theta)=\lambda f'(\lambda),
\qquad
g''(\theta)=\lambda f'(\lambda)+\lambda^2f''(\lambda).
\]
This preserves positivity and makes multiplicative penalty scales easier to
search.

\section{Multiple stationary points}
A validation objective need not be unimodal. The active numerical strategy is
therefore to use a coarse log-penalty grid to locate derivative sign changes,
apply a bracketed root solver to each interval, classify the stationary points,
and compare all local minima with the search boundaries. Newton and direct
bounded minimization remain useful comparison methods.

\section{Chronological validation}
Rolling-origin evaluation repeats
\[
1{:}T_j
\longrightarrow
T_j{+}1{:}T_j{+}h
\]
at several origins. This separates genuine forecasting from interpolation or
trend reconstruction with future observations available on both sides.

\section{Implementation map}
Reusable mathematics belongs in \texttt{src/trend\_estimation}. Internal
derivations and research decisions live in \texttt{notes/}. The active
research manuscript lives in \texttt{paper\_forecast-optimal-smoothing/}.

\end{document}
